\documentclass[11pt]{article}
\usepackage[margin=1in]{geometry}
\usepackage{amsmath,amssymb,amsthm,hyperref}
\newtheorem{theorem}{Theorem}
\newtheorem{lemma}{Lemma}
\newcommand{\conv}{\operatorname{conv}}
\newcommand{\W}{\mathcal W}
\title{Even Wachs permutation polytopes are products}
\author{Principia Math In Progress}
\date{27 September 2026}
\begin{document}
\maketitle
\begin{abstract}
The convex hull of the Wachs permutations in $S_{2m}$ is affinely
isomorphic to the product of the ordinary $(m-1)$-dimensional
permutahedron and an $m$-cube. In particular it is simple for every
$m\geq1$, proving Conjecture 6.3 of Brenti and Sentinelli.
\end{abstract}

\section*{Task 1: the source conjecture}
For $i\in[2m]$, put $i^*=i+1$ if $i$ is odd and $i^*=i-1$ if $i$ is even.
The source definition is
\[
 \W(S_{2m})=\{w\in S_{2m}:
 |w^{-1}(i)-w^{-1}(i^*)|\leq1\text{ for every }i\in[2m-1]\}.
\]
We identify $w$ with $(w(1),\ldots,w(2m))\in\mathbb R^{2m}$ and set
$P_m=\conv\W(S_{2m})$. Conjecture 6.3 of \cite{BS} asks whether
$P_m$ is simple for every positive integer $m$. Simplicity is measured
in the affine hull: a $d$-dimensional polytope is simple if each vertex
is incident with exactly $d$ edges.

\begin{theorem}\label{main}
Let
\[
 \Pi_m=\conv\{(\sigma(1),\ldots,\sigma(m)):\sigma\in S_m\},
 \qquad C_m=[-1,1]^m.
\]
Then $P_m$ is affinely isomorphic to $\Pi_m\times C_m$. It has dimension
$2m-1$, has $2^m m!$ vertices, and is simple. Consequently its number
of edges is $(2m-1)2^{m-1}m!$.
\end{theorem}

\begin{lemma}\label{blocks}
Every $w\in\W(S_{2m})$ is uniquely obtained by arranging the $m$ pairs
$\{2j-1,2j\}$ in an arbitrary order and choosing the order inside each
pair independently.
\end{lemma}
\begin{proof}
The positions occupied by $2j-1$ and $2j$ are distinct and have distance
at most one, so they are consecutive. These disjoint position pairs cover
$[2m]$. Position 1 must be paired with position 2; removing them and
repeating shows that the position pairs are precisely
$\{1,2\},\{3,4\},\ldots,\{2m-1,2m\}$. The value pair placed in each
position pair is arbitrary, as is its orientation. Conversely, every such
arrangement satisfies every condition in the definition.
\end{proof}

\section*{The affine product}
For $x\in\mathbb R^{2m}$ define $T(x)=(y,z)\in\mathbb R^m\times\mathbb R^m$
by
\begin{equation}\label{map}
 y_i=\frac{x_{2i-1}+x_{2i}+1}{4},\qquad
 z_i=x_{2i}-x_{2i-1}\quad(1\leq i\leq m).
\end{equation}
This affine map is invertible, with
\begin{equation}\label{inverse}
 x_{2i-1}=2y_i-\frac12-\frac{z_i}{2},\qquad
 x_{2i}=2y_i-\frac12+\frac{z_i}{2}.
\end{equation}
By Lemma \ref{blocks}, on the Wachs set the vector $y$ runs over all
permutations of $(1,\ldots,m)$ and $z$ runs independently over
$\{-1,1\}^m$. Hence
\[
 T(\W(S_{2m}))=S_m\times\{-1,1\}^m,
\]
where $S_m$ on the right denotes its one-line coordinate vectors.
For finite sets $A,B$ one has
$\conv(A\times B)=\conv(A)\times\conv(B)$: the nontrivial inclusion
follows by multiplying convex-combination weights,
$(\sum_a\lambda_a a,\sum_b\mu_b b)
=\sum_{a,b}\lambda_a\mu_b(a,b)$.
Affine maps preserve convex combinations, proving
\begin{equation}\label{product}
 T(P_m)=\Pi_m\times C_m.
\end{equation}

\section*{Dimension and simplicity, with an explicit cone proof}
Write $v=(1,\ldots,m)$ and $g_k=e_k-e_{k+1}$ for $1\leq k<m$.
For a permutation vector $u$, define
\[
 a_k=\sum_{i=1}^k(u_i-i)\geq0.
\]
The inequality holds because any $k$ distinct entries of $[m]$ have
sum at least $1+\cdots+k$. Since $\sum_i(u_i-i)=0$, telescoping gives
\begin{equation}\label{cone}
 u-v=\sum_{k=1}^{m-1}a_k g_k.
\end{equation}
Each $g_k$ itself is a vertex difference, obtained by interchanging the
entries $k,k+1$ in $v$. These $m-1$ vectors are linearly independent.
Thus the cone of feasible directions of $\Pi_m$ at $v$ is exactly their
nonnegative span and has dimension $m-1$. Coordinate permutations act
transitively on the permutation vectors. Each vector is a vertex: the
functional $u\mapsto v\cdot u$ is uniquely maximized at $v$, since
\[
 v\cdot u=\|v\|^2-\tfrac12\|u-v\|^2
\]
for permutation vectors $u$, and the same argument applies after
permuting coordinates.

At the vertex $(v,\mathbf1)$ of the product in \eqref{product}, the cone
of feasible directions is generated by
\[
 (g_k,0)\ (1\leq k<m),\qquad (0,-2e_i)\ (1\leq i\leq m).
\]
Indeed, \eqref{cone} supplies the first components, while each cube vertex
$z$ satisfies $z-\mathbf1=\sum_i(1-z_i)(-2e_i)/2$.
Every displayed generator is itself a vertex difference. The $2m-1$
generators are linearly independent, so this cone has exactly $2m-1$
extreme rays. Extreme rays of a polytope's feasible-direction cone at
a vertex correspond to its incident edges: an exposing functional for
such a ray cuts out the corresponding one-dimensional face, and the
converse follows by exposing an edge.
Consequently the product has dimension $2m-1$ and exactly $2m-1$ edges
at this vertex. Coordinate permutations in the first factor and sign
changes in the second act transitively on all its vertices, proving
simplicity everywhere. The vertices of a product are exactly the pairs
of vertices of its factors, giving $2^m m!$ of them. Counting incidences
of vertices and edges proves the stated edge count. This also covers
$m=1$, when $\Pi_1$ is a point and $P_1$ is a segment, and completes
the proof of Theorem \ref{main}.

\section*{Scope and verification}
The argument uses the even-index definition and proves the full positive
integer range. It makes no assertion about the odd-index hulls or the
Bruhat order. The independent finite verifier exhausts $S_{2m}$ for
$1\leq m\leq4$ using the original inverse-position condition, checks
the product identity and inverse map, and checks the exact cone
certificates through $m=5$. These finite checks support the universal
proof above; they do not replace it. This is author-agent review,
without independent peer review or a proof assistant.

\section*{Proposed follow-up tasks}
\begin{description}
\item[Task 2.] Give a closed formula for every face number of $P_m$ by
combining ordered set partitions for $\Pi_m$ with the faces of $C_m$.
\item[Task 3.] Determine the Ehrhart polynomial of $P_m$ for the lattice
$\mathbb Z^{2m}\cap\operatorname{aff}(P_m)$, accounting for the lattice
change under \eqref{map}.
\item[Task 4.] For real numbers $a_1<b_1<a_2<b_2<\cdots<a_m<b_m$,
classify when the hull of all ordered, independently oriented blocks
$(a_j,b_j)$ is simple, including the case of unequal block lengths.
\item[Task 5.] Determine exactly which odd integers $n$ have simple
$\conv\W(S_n)$; for every remaining $n$, exhibit a vertex with more
incident edges than the affine dimension.
\end{description}

\begin{thebibliography}{9}
\bibitem{BS} F. Brenti and P. Sentinelli,
\emph{Wachs permutations, Bruhat order and weak order},
European Journal of Combinatorics \textbf{119} (2024), 103804,
\href{https://doi.org/10.1016/j.ejc.2023.103804}{doi:10.1016/j.ejc.2023.103804}.
\href{https://arxiv.org/abs/2212.04932}{arXiv:2212.04932v1};
Conjecture 6.3, p.~36. The Polimi repository supplies a byte-identical
38-page file labeled postprint/author's accepted manuscript;
its version provenance and the separately read 31-page author preprint
are detailed in the accompanying literature audit.
\end{thebibliography}
\end{document}
